\section{Identificación y definición nominal-conceptual de variables}

\subsection{Tiempo de exposición a pantallas digitales}

\textbf{Definición nominal:}

Cantidad de tiempo durante el cual los estudiantes utilizan dispositivos electrónicos con pantalla en sus actividades académicas y personales.

\textbf{Definición conceptual:}

Hace referencia al número de horas dedicadas al uso de computadores, teléfonos inteligentes u otros dispositivos digitales. La exposición prolongada a pantallas ha sido relacionada con síntomas visuales y molestias derivadas del uso continuo de tecnología digital [2], [3].


\subsection{Hábitos posturales durante la jornada académica}

\textbf{Definición nominal:}

Forma en que los estudiantes mantienen su postura corporal durante actividades académicas realizadas frente a dispositivos electrónicos.

\textbf{Definición conceptual:}

Comprende las posiciones adoptadas por los estudiantes durante periodos prolongados de estudio, incluyendo postura de espalda, cuello, brazos y muñecas. Las condiciones posturales inadecuadas pueden relacionarse con molestias musculoesqueléticas [1].


\subsection{Síntomas físicos asociados al uso prolongado de pantallas}

\textbf{Definición nominal:}

Manifestaciones físicas experimentadas por los estudiantes relacionadas con el uso frecuente de dispositivos digitales.

\textbf{Definición conceptual:}

Incluye síntomas como fatiga visual, irritación ocular, dolor cervical, molestias en hombros, espalda y extremidades superiores asociados a jornadas prolongadas frente a pantallas [2], [3].
\estado{En construcción}
\notafase{Se redacta en la Fase 3 del curso.}
